\documentclass[11pt]{article}

\usepackage[a4paper,margin=1in]{geometry}
\usepackage{amsmath,amssymb,amsthm,mathtools}
\usepackage{bm}
\usepackage{mathrsfs}
\usepackage{hyperref}
\usepackage{enumitem}
\usepackage{booktabs}
\usepackage{microtype}

\hypersetup{
  colorlinks=true,
  linkcolor=blue,
  citecolor=blue,
  urlcolor=blue
}

\newtheorem{theorem}{Theorem}
\newtheorem{proposition}[theorem]{Proposition}
\newtheorem{lemma}[theorem]{Lemma}
\newtheorem{corollary}[theorem]{Corollary}
\theoremstyle{definition}
\newtheorem{definition}[theorem]{Definition}
\newtheorem{assumption}[theorem]{Assumption}
\newtheorem{remark}[theorem]{Remark}

\newcommand{\R}{\mathbb{R}}
\newcommand{\E}{\mathbb{E}}
\newcommand{\Prob}{\mathbb{P}}
\newcommand{\supp}{\mathrm{supp}}
\newcommand{\norm}[1]{\left\lVert #1 \right\rVert}
\newcommand{\mc}[1]{\mathcal{#1}}
\newcommand{\cE}{\mathcal{E}}
\newcommand{\cB}{\mathcal{B}}

\title{Linear Mode Connectivity for the Single-Gaussian\\Natural-Softplus Head}
\author{Working draft}
\date{\today}

\begin{document}
\maketitle

\begin{abstract}
This note isolates the $K=1$ part of the linear-mode-connectivity program for
distributional two-layer networks with Gaussian natural outputs. The training
side is assumed to be supplied by the companion mean-field note: in the
continuous-time full-batch regime, two wide runs started from the same
initialization law track the same deterministic path. The present note studies
the interpolation side. For the practical single-Gaussian head
\[
  \eta_1(x)=\frac1N\sum_{i=1}^N a_i\varphi(\xi_i,x),
  \qquad
  \eta_2(x)=-\psi_{\rm sp}\!\left(\frac1N\sum_{i=1}^N c_i\varphi(\xi_i,x)\right),
\]
the Gaussian negative log-likelihood is convex in the natural output
$\eta=(\eta_1,\eta_2)$. The barrier
instead decomposes into two terms: a vector-head transfer gap between parameter
interpolation and raw-output interpolation, and a softplus Jensen gap between
raw-output interpolation and natural-output interpolation. Combining this
deterministic decomposition with the companion common-path theorem yields a
Ferbach-style conclusion: if two trained width-$N$ networks come from the same
initialization law, then after an optimal neuron permutation their
interpolation barrier goes to zero as $N\to\infty$.
\end{abstract}

\section{Executive Summary}

For the single-Gaussian head, the LMC story is cleaner than in the multi-mixture
case:
\begin{enumerate}[label=(\roman*)]
  \item there is no component permutation and no softmax competition;
  \item the Gaussian negative log-likelihood is convex in the natural output
  $\eta=(\eta_1,\eta_2)$;
  \item the practical parameterization still uses a raw scale field
  $s$ with $\eta_2=-\psi_{\rm sp}(s)$, so the parameter-interpolation path is
  not the same as the natural-output interpolation path;
  \item the interpolation barrier therefore splits into exactly two terms:
  a vector-head transfer term and a softplus Jensen term;
  \item the companion mean-field theorem supplies the missing training-side
  input, namely that two runs with the same initialization law are close in
  Wasserstein distance at any fixed finite horizon.
\end{enumerate}

The main message is:
\begin{quote}
In the $K=1$ Gaussian natural-softplus model, one can prove a mode-connectivity
theorem directly in
natural-output space; all that remains is to quantify the mismatch between
parameter interpolation and natural-output interpolation.
\end{quote}

\section{Setup}

Let $\mc{X}$ be the input space and let $(X,Y)$ be a data pair with
$|Y|\le Y_\ast$ almost surely. Fix a bounded feature map
$\varphi:\R^p\times\mc{X}\to\R$ satisfying
\[
  K_\varphi:=\sup_{\xi,x}|\varphi(\xi,x)|<\infty,
  \qquad
  K_{\nabla}:=\sup_{\xi,x}\norm{\nabla_\xi\varphi(\xi,x)}<\infty.
\]
We consider width-$N$ two-layer networks with mean-field scaling
\[
  z_\Theta(x)
  =
  \bigl(u_\Theta(x),s_\Theta(x)\bigr)
  :=
  \frac1N\sum_{i=1}^N (a_i,c_i)\,\varphi(\xi_i,x)
  \in \R^2,
\]
where
\[
  \Theta=((a_i,c_i,\xi_i))_{i=1}^N \in (\R\times\R\times\R^p)^N.
\]
The natural output is obtained by the practical softplus link
\[
  \eta_\Theta(x)
  :=
  \Gamma(z_\Theta(x))
  =
  \bigl(u_\Theta(x),-\psi_{\rm sp}(s_\Theta(x))\bigr),
\]
where
\[
  \psi_{\rm sp}(s):=\lambda_{\min}+\log(1+e^s),
  \qquad
  \lambda_{\min}>0.
\]

\paragraph{Population risk.}
For Gaussian natural parameters
$\eta=(\eta_1,\eta_2)\in \R\times(-\infty,0)$, define
\[
  A_{\rm G}(\eta_1,\eta_2)
  =
  -\frac{\eta_1^2}{4\eta_2}
  + \frac12\log\!\Bigl(\frac{-\pi}{\eta_2}\Bigr),
\]
and the negative log-likelihood
\[
  \ell_\eta(\eta;y)
  :=
  A_{\rm G}(\eta) - \eta_1 y - \eta_2 y^2.
\]
The population risk is
\[
  \cE(\Theta)
  :=
  \E\big[\ell_\eta(\eta_\Theta(X);Y)\big].
\]

\paragraph{Aligned interpolation.}
Fix a neuron permutation $\pi\in\mathfrak S_N$ and absorb it into network $B$;
that is, after relabeling we write
\[
  \Theta_A=((a_i^A,c_i^A,\xi_i^A))_{i=1}^N,
  \qquad
  \Theta_B=((a_i^B,c_i^B,\xi_i^B))_{i=1}^N.
\]
The parameter-interpolation path is
\[
  \Theta_t
  :=
  \bigl((1-t)a_i^A+ta_i^B,\ (1-t)c_i^A+tc_i^B,\ (1-t)\xi_i^A+t\xi_i^B\bigr)_{i=1}^N.
\]
We also define
\[
  z_A:=z_{\Theta_A},
  \qquad
  z_B:=z_{\Theta_B},
  \qquad
  z_t:=z_{\Theta_t},
\]
the raw-output interpolation
\[
  \hat z_t := (1-t)z_A + tz_B,
\]
the natural-output interpolation
\[
  \hat\eta_t := (1-t)\eta_{\Theta_A} + t\eta_{\Theta_B},
\]
and the lifted raw interpolation
\[
  \tilde\eta_t := \Gamma(\hat z_t).
\]

\begin{definition}[Aligned barrier]
For a fixed neuron permutation $\pi$, the aligned interpolation barrier is
\[
  \cB(\Theta_A,\Theta_B;\pi)
  :=
  \sup_{t\in[0,1]}
  \Bigl\{
    \cE(\Theta_t)
    -
    (1-t)\cE(\Theta_A)
    -
    t\cE(\Theta_B)
  \Bigr\}.
\]
\end{definition}

\paragraph{Alignment diagnostics.}
The hidden-feature mismatch is
\begin{equation}
  B_N
  :=
  \frac1N\sum_{i=1}^N \norm{\xi_i^A-\xi_i^B}^2.
  \label{eq:hidden_mismatch}
\end{equation}
The raw-scale endpoint gap is
\begin{equation}
  \Delta_{s,N}
  :=
  \E\bigl[(s_{\Theta_A}(X)-s_{\Theta_B}(X))^2\bigr].
  \label{eq:scale_gap}
\end{equation}
Finally, define the per-coordinate head moment
\begin{equation}
  M_V^\infty
  :=
  \max\!\Biggl\{
    \Bigl(\frac1N\sum_{i=1}^N (a_i^A)^2\Bigr)^{1/2},
    \Bigl(\frac1N\sum_{i=1}^N (c_i^A)^2\Bigr)^{1/2},
    \Bigl(\frac1N\sum_{i=1}^N (a_i^B)^2\Bigr)^{1/2},
    \Bigl(\frac1N\sum_{i=1}^N (c_i^B)^2\Bigr)^{1/2}
  \Biggr\}.
  \label{eq:head_moment}
\end{equation}

\section{Output Geometry for the Single-Gaussian Head}

\begin{proposition}[Convexity and Lipschitz control on a natural strip]
\label{prop:convex-lipschitz-strip}
Fix $M>0$ and $\lambda_{\min}>0$, and define
\[
  D_{M,\lambda_{\min}}
  :=
  \{(\eta_1,\eta_2)\in\R^2: |\eta_1|\le M,\ \eta_2\le -\lambda_{\min}\}.
\]
Then for every $|y|\le Y_\ast$:
\begin{enumerate}[label=(\roman*)]
  \item $\ell_\eta(\cdot;y)$ is convex on $D_{M,\lambda_{\min}}$;
  \item $\ell_\eta(\cdot;y)$ is $C_\eta(M,\lambda_{\min},Y_\ast)$-Lipschitz in
  the $\ell_\infty$ norm on $D_{M,\lambda_{\min}}$, where
  \begin{align}
    C_\eta(M,\lambda_{\min},Y_\ast)
    &:=
    \frac{M}{2\lambda_{\min}} + Y_\ast
    + \frac{M^2}{4\lambda_{\min}^2}
    + \frac{1}{2\lambda_{\min}}
    + Y_\ast^2.
    \label{eq:Ceta}
  \end{align}
\end{enumerate}
\end{proposition}

\begin{proof}
For Gaussian natural parameters, $\ell_\eta(\eta;y)=A_{\rm G}(\eta)-\eta_1y-\eta_2y^2$.
The Gaussian log-partition $A_{\rm G}$ is convex on $\{\eta_2<0\}$, so
$\ell_\eta(\cdot;y)$ is convex there as well.

For the Lipschitz bound, write
\[
  q_1(\eta;y):=\partial_{\eta_1}\ell_\eta(\eta;y)
  =
  -\frac{\eta_1}{2\eta_2}-y
\]
and
\[
  q_2(\eta;y):=\partial_{\eta_2}\ell_\eta(\eta;y)
  =
  \frac{\eta_1^2}{4\eta_2^2}
  -
  \frac1{2\eta_2}
  -
  y^2.
\]
On $D_{M,\lambda_{\min}}$ and for $|y|\le Y_\ast$,
\[
  |q_1(\eta;y)|
  \le
  \frac{M}{2\lambda_{\min}} + Y_\ast
\]
and
\[
  |q_2(\eta;y)|
  \le
  \frac{M^2}{4\lambda_{\min}^2}
  +
  \frac1{2\lambda_{\min}}
  +
  Y_\ast^2.
\]
Hence
\[
  \sup_{\eta\in D_{M,\lambda_{\min}}}
  \norm{\nabla_\eta \ell_\eta(\eta;y)}_1
  \le
  C_\eta(M,\lambda_{\min},Y_\ast).
\]
The mean-value theorem in the $\ell_\infty$ norm yields the claim.
\end{proof}

\begin{lemma}[Automatic strip control along the interpolation path]
\label{lem:auto_strip}
For all $x\in\mc{X}$ and all $t\in[0,1]$,
\[
  \norm{z_A(x)}_\infty,\ \norm{z_B(x)}_\infty,\ \norm{z_t(x)}_\infty,\
  \norm{\hat z_t(x)}_\infty
  \le
  K_\varphi M_V^\infty.
\]
Consequently
\[
  \eta_{\Theta_A}(x),\ \eta_{\Theta_B}(x),\ \eta_{\Theta_t}(x),\
  \hat\eta_t(x),\ \tilde\eta_t(x)
  \in
  D_{K_\varphi M_V^\infty,\lambda_{\min}}
  \qquad
  \forall x,\ \forall t.
\]
\end{lemma}

\begin{proof}
Consider the first coordinate. By Cauchy--Schwarz,
\[
  |u_{\Theta_A}(x)|
  \le
  \frac1N\sum_{i=1}^N |a_i^A|\,|\varphi(\xi_i^A,x)|
  \le
  K_\varphi
  \Bigl(\frac1N\sum_{i=1}^N (a_i^A)^2\Bigr)^{1/2}
  \le
  K_\varphi M_V^\infty.
\]
The same estimate holds for the raw scale coordinate and for network $B$.
For the interpolated head coefficients, convexity of the square gives
\[
  \frac1N\sum_{i=1}^N \bigl((1-t)a_i^A+ta_i^B\bigr)^2
  \le
  (1-t)\frac1N\sum_{i=1}^N (a_i^A)^2
  +
  t\frac1N\sum_{i=1}^N (a_i^B)^2
  \le
  (M_V^\infty)^2,
\]
and similarly for the $c$-coordinates. This proves the raw bound for $z_t$.
The same bound for $\hat z_t=(1-t)z_A+tz_B$ is immediate.

For the natural outputs, the first coordinate is unchanged and the second
coordinate is always $\le -\lambda_{\min}$ because
$\psi_{\rm sp}(s)\ge \lambda_{\min}$ for all $s\in\R$.
\end{proof}

\begin{proposition}[Softplus Jensen gap]
\label{prop:softplus-jensen}
For any $s_A,s_B\in\R$ and $t\in[0,1]$,
\begin{equation}
  0
  \le
  (1-t)\psi_{\rm sp}(s_A)+t\psi_{\rm sp}(s_B)
  -
  \psi_{\rm sp}((1-t)s_A+ts_B)
  \le
  \frac18\,t(1-t)(s_A-s_B)^2.
  \label{eq:softplus-jensen}
\end{equation}
Consequently, for every $x\in\mc{X}$,
\begin{equation}
  \norm{\tilde\eta_t(x)-\hat\eta_t(x)}_\infty
  \le
  \frac18\,t(1-t)\bigl(s_{\Theta_A}(x)-s_{\Theta_B}(x)\bigr)^2.
  \label{eq:natural-jensen-gap}
\end{equation}
\end{proposition}

\begin{proof}
The function $\psi_{\rm sp}$ is convex and satisfies
\[
  \psi_{\rm sp}''(s)
  =
  \frac{e^s}{(1+e^s)^2}
  \le
  \frac14.
\]
For any $C^2$ convex function with second derivative bounded by $L$,
the midpoint interpolation error is at most
$\frac{L}{2}t(1-t)(s_A-s_B)^2$. Applying this with
$L=\frac14$ gives \eqref{eq:softplus-jensen}.
Since only the second natural coordinate is nonlinear,
\eqref{eq:natural-jensen-gap} is the same estimate written for
\[
  \tilde\eta_t
  =
  \bigl((1-t)u_A+tu_B,\ -\psi_{\rm sp}((1-t)s_A+ts_B)\bigr)
\]
and
\[
  \hat\eta_t
  =
  \bigl((1-t)u_A+tu_B,\ -(1-t)\psi_{\rm sp}(s_A)-t\psi_{\rm sp}(s_B)\bigr).
\]
\end{proof}

\section{Vector-Head Transfer}

The raw head is two-dimensional: one coordinate feeds the first natural
parameter directly, while the other feeds the raw scale. This is exactly the
vector-head setting in which the direct output-coupling argument is useful.

\begin{proposition}[Direct raw-output coupling for the two-coordinate head]
\label{prop:vector-head-transfer}
For the aligned pair $(\Theta_A,\Theta_B)$,
\begin{equation}
  \sup_{t\in[0,1]}\ \sup_{x\in\mc{X}}
  \norm{z_t(x)-\hat z_t(x)}_\infty
  \le
  \frac12 K_{\nabla} M_V^\infty \sqrt{B_N}.
  \label{eq:vector-transfer}
\end{equation}
In particular,
\[
  \sup_{t\in[0,1]}
  \Bigl(
    \E\big[\norm{z_t(X)-\hat z_t(X)}_\infty^2\big]
  \Bigr)^{1/2}
  \le
  \frac12 K_{\nabla} M_V^\infty \sqrt{B_N}.
\]
\end{proposition}

\begin{proof}
Set
\[
  \xi_i(t):=(1-t)\xi_i^A+t\xi_i^B,
  \qquad
  \Delta \xi_i:=\xi_i^A-\xi_i^B.
\]
Let $v_i^A=(a_i^A,c_i^A)$ and $v_i^B=(a_i^B,c_i^B)$. Then
\begin{align*}
  z_t(x)-\hat z_t(x)
  &=
  \frac1N\sum_{i=1}^N
  \Bigl[
    (1-t)v_i^A\bigl(\varphi(\xi_i(t),x)-\varphi(\xi_i^A,x)\bigr)
    \\
  &\qquad\qquad\qquad
    +
    tv_i^B\bigl(\varphi(\xi_i(t),x)-\varphi(\xi_i^B,x)\bigr)
  \Bigr].
\end{align*}
Since
\[
  \norm{\xi_i(t)-\xi_i^A}=t\norm{\Delta\xi_i},
  \qquad
  \norm{\xi_i(t)-\xi_i^B}=(1-t)\norm{\Delta\xi_i},
\]
the uniform gradient bound gives
\[
  |\varphi(\xi_i(t),x)-\varphi(\xi_i^A,x)|
  \le
  tK_{\nabla}\norm{\Delta\xi_i},
  \qquad
  |\varphi(\xi_i(t),x)-\varphi(\xi_i^B,x)|
  \le
  (1-t)K_{\nabla}\norm{\Delta\xi_i}.
\]
Taking the $j$th coordinate ($j\in\{1,2\}$), we obtain
\begin{align*}
  |(z_t(x)-\hat z_t(x))_j|
  &\le
  \frac{t(1-t)K_{\nabla}}{N}
  \sum_{i=1}^N
  \bigl(|(v_i^A)_j|+|(v_i^B)_j|\bigr)\norm{\Delta\xi_i}.
\end{align*}
By Cauchy--Schwarz,
\[
  \frac1N\sum_{i=1}^N |(v_i^A)_j|\norm{\Delta\xi_i}
  \le
  \Bigl(\frac1N\sum_{i=1}^N (v_i^A)_j^2\Bigr)^{1/2}
  \Bigl(\frac1N\sum_{i=1}^N \norm{\Delta\xi_i}^2\Bigr)^{1/2}
  \le
  M_V^\infty \sqrt{B_N},
\]
and the same bound holds for the $B$ term. Hence
\[
  \norm{z_t(x)-\hat z_t(x)}_\infty
  \le
  2t(1-t)K_{\nabla}M_V^\infty\sqrt{B_N}
  \le
  \frac12 K_{\nabla}M_V^\infty\sqrt{B_N}.
\]
The $L^2(P_X)$ statement follows immediately.
\end{proof}

\begin{proposition}[Natural-output interpolation gap]
\label{prop:natural-output-gap}
For every $t\in[0,1]$,
\begin{equation}
  \E\big[\norm{\eta_{\Theta_t}(X)-\hat\eta_t(X)}_\infty\big]
  \le
  \frac12 K_{\nabla} M_V^\infty \sqrt{B_N}
  +
  \frac18\,t(1-t)\Delta_{s,N}.
  \label{eq:natural-gap}
\end{equation}
\end{proposition}

\begin{proof}
By the triangle inequality,
\[
  \norm{\eta_{\Theta_t}(x)-\hat\eta_t(x)}_\infty
  \le
  \norm{\eta_{\Theta_t}(x)-\tilde\eta_t(x)}_\infty
  +
  \norm{\tilde\eta_t(x)-\hat\eta_t(x)}_\infty.
\]
The map $\Gamma(u,s)=(u,-\psi_{\rm sp}(s))$ is $1$-Lipschitz in the
$\ell_\infty$ norm because its first coordinate is the identity and
$0<\psi_{\rm sp}'(s)<1$. Hence
\[
  \norm{\eta_{\Theta_t}(x)-\tilde\eta_t(x)}_\infty
  \le
  \norm{z_t(x)-\hat z_t(x)}_\infty,
\]
so Proposition~\ref{prop:vector-head-transfer} gives
\[
  \E\big[\norm{\eta_{\Theta_t}(X)-\tilde\eta_t(X)}_\infty\big]
  \le
  \frac12 K_{\nabla} M_V^\infty \sqrt{B_N}.
\]
For the second term, Proposition~\ref{prop:softplus-jensen} yields
\[
  \norm{\tilde\eta_t(x)-\hat\eta_t(x)}_\infty
  \le
  \frac18\,t(1-t)(s_{\Theta_A}(x)-s_{\Theta_B}(x))^2.
\]
Taking expectation proves \eqref{eq:natural-gap}.
\end{proof}

\section{Deterministic LMC Theorem}

\begin{theorem}[Aligned barrier decomposition for the single-Gaussian head]
\label{thm:deterministic-lmc}
For every fixed neuron permutation $\pi$,
\begin{equation}
  \cB(\Theta_A,\Theta_B;\pi)
  \le
  C_\eta(K_\varphi M_V^\infty,\lambda_{\min},Y_\ast)
  \left[
    \frac12 K_{\nabla} M_V^\infty \sqrt{B_N}
    +
    \frac1{32}\Delta_{s,N}
  \right].
  \label{eq:deterministic-barrier}
\end{equation}
\end{theorem}

\begin{proof}
Fix $t\in[0,1]$. By Lemma~\ref{lem:auto_strip}, all relevant outputs lie in
$D_{K_\varphi M_V^\infty,\lambda_{\min}}$, so
Proposition~\ref{prop:convex-lipschitz-strip} applies there.
Using the Lipschitz bound first,
\[
  \E[\ell_\eta(\eta_{\Theta_t}(X);Y)]
  \le
  \E[\ell_\eta(\hat\eta_t(X);Y)]
  +
  C_\eta(K_\varphi M_V^\infty,\lambda_{\min},Y_\ast)
  \E[\norm{\eta_{\Theta_t}(X)-\hat\eta_t(X)}_\infty].
\]
Proposition~\ref{prop:natural-output-gap} bounds the second term by
\[
  C_\eta(K_\varphi M_V^\infty,\lambda_{\min},Y_\ast)
  \left[
    \frac12 K_{\nabla} M_V^\infty \sqrt{B_N}
    +
    \frac18\,t(1-t)\Delta_{s,N}
  \right].
\]
Now use convexity of $\ell_\eta(\cdot;Y)$ in natural coordinates:
\[
  \ell_\eta(\hat\eta_t(X);Y)
  \le
  (1-t)\ell_\eta(\eta_{\Theta_A}(X);Y)
  +
  t\ell_\eta(\eta_{\Theta_B}(X);Y).
\]
Taking expectations and then $\sup_{t\in[0,1]}$ gives
\eqref{eq:deterministic-barrier} because $\sup_{t\in[0,1]} t(1-t)=1/4$.
\end{proof}

\begin{remark}[Barrier ingredients]
Theorem~\ref{thm:deterministic-lmc} is the $K=1$ analogue of the
Ferbach-style barrier decomposition. Its two terms have a direct interpretation:
\begin{enumerate}[label=(\alph*)]
  \item convexity is exact in natural-output space, so no output-level
  curvature loss appears;
  \item the only nonlinear residue is the softplus Jensen term
  $\Delta_{s,N}$;
  \item the head is genuinely vector-valued, and the transfer step therefore
  uses the direct vector-head coupling proposition rather than any scalar-head
  activation-matching argument.
\end{enumerate}
\end{remark}

\section{Rates}

The deterministic theorem is expressed in terms of the alignment quantities
$B_N$ and $\Delta_{s,N}$. The width rate therefore depends on how those two
objects scale along the aligned pair.

\begin{proposition}[From hidden-law concentration to $B_N=O(d/N)$]
\label{prop:BN-from-w2}
Let
\[
  \hat\nu_A^N:=\frac1N\sum_{i=1}^N \delta_{\xi_i^A},
  \qquad
  \hat\nu_B^N:=\frac1N\sum_{i=1}^N \delta_{\xi_i^B},
\]
and let $\mu_{\xi,\ast}$ be a reference law on $\R^p$. Suppose that for some
deterministic sequence $\eta_N\downarrow 0$,
\[
  W_2(\hat\nu_A^N,\mu_{\xi,\ast})\le \eta_N,
  \qquad
  W_2(\hat\nu_B^N,\mu_{\xi,\ast})\le \eta_N.
\]
If the neuron permutation $\pi_N$ is chosen by hidden-weight optimal transport,
then
\begin{equation}
  B_N = W_2^2(\hat\nu_A^N,\hat\nu_B^N)\le 4\eta_N^2.
  \label{eq:BN-via-w2}
\end{equation}
In particular, if
\[
  \eta_N = O(\kappa_\xi \sqrt{d/N}),
\]
then
\[
  B_N = O(\kappa_\xi^2 d/N).
\]
\end{proposition}

\begin{proof}
Because the hidden states all carry mass $1/N$, the $W_2$-optimal coupling
between $\hat\nu_A^N$ and $\hat\nu_B^N$ is realized by a permutation, and the
corresponding squared cost is exactly $B_N$. The triangle inequality gives
\[
  W_2(\hat\nu_A^N,\hat\nu_B^N)
  \le
  W_2(\hat\nu_A^N,\mu_{\xi,\ast}) + W_2(\hat\nu_B^N,\mu_{\xi,\ast})
  \le
  2\eta_N.
\]
Squaring proves \eqref{eq:BN-via-w2}. The final claim is immediate.
\end{proof}

\begin{corollary}[Baseline rate]
\label{cor:baseline-rate}
Suppose that along a width sequence,
\[
  B_N = O(d/N),
  \qquad
  \Delta_{s,N} = O(1/N).
\]
Then Theorem~\ref{thm:deterministic-lmc} yields
\[
  \cB(\Theta_A,\Theta_B;\pi)
  =
  O\!\left(\sqrt{\frac{d}{N}} + \frac1N\right).
\]
In particular, for fixed input dimension $d$,
\[
  \cB(\Theta_A,\Theta_B;\pi)=O(N^{-1/2}).
\]
\end{corollary}

\begin{proof}
Substitute the displayed assumptions into
\eqref{eq:deterministic-barrier}.
\end{proof}

\begin{remark}[Ferbach-style specialization]
Proposition~\ref{prop:BN-from-w2} is the precise point where the present note
meets the Ferbach-style alignment theory. If the trained hidden empirical laws
obey a sub-Gaussian Wasserstein concentration rate
\[
  \eta_N=O(\kappa_\xi\sqrt{d/N}),
\]
then hidden-weight optimal transport gives
\[
  B_N = O(\kappa_\xi^2 d/N).
\]
Substituting this into Corollary~\ref{cor:baseline-rate} yields the explicit
baseline barrier rate
\[
  \cB(\Theta_A,\Theta_B;\pi_N)
  =
  O\!\left(\kappa_\xi\sqrt{\frac{d}{N}} + \frac1N\right).
\]
Thus the $O(\sqrt{d/N})$ law comes entirely from the transfer bottleneck.
\end{remark}

The older LMC draft identified a faster route: the transfer term can improve
from $O(\sqrt{B_N})$ to $O(B_N)$ if the trained two-coordinate head is itself a
regular function of the hidden feature. For $K=1$, this route is simpler than
in the mixture case because there is no softmax block and hence no gauge-fixing
step.

\begin{definition}[Per-coordinate head mismatch]
\label{def:head-mismatch}
Define
\[
  M_{\Delta V}^\infty
  :=
  \max_{j\in\{1,2\}}
  \Bigl(
    \frac1N\sum_{i=1}^N
    \bigl((v_i^A-v_i^B)_j\bigr)^2
  \Bigr)^{1/2},
  \qquad
  v_i^X:=(a_i^X,c_i^X)\in\R^2.
\]
\end{definition}

\begin{assumption}[Conditional head regularity for $K=1$]
\label{ass:head-regularity-k1}
There exist a map $g:\R^p\to\R^2$ and residuals $r_i^A,r_i^B\in\R^2$ such that
\[
  v_i^A = g(\xi_i^A)+r_i^A,
  \qquad
  v_i^B = g(\xi_i^B)+r_i^B,
\]
and
\[
  \norm{g(\xi)-g(\bar\xi)}_\infty \le L_g \norm{\xi-\bar\xi}
  \qquad
  \forall \xi,\bar\xi\in\R^p.
\]
Define the residual mismatch
\[
  \mathcal R_N
  :=
  \max_{j\in\{1,2\}}
  \Bigl(
    \frac1N\sum_{i=1}^N
    \bigl((r_i^A-r_i^B)_j\bigr)^2
  \Bigr)^{1/2}.
\]
\end{assumption}

\begin{lemma}[Head mismatch from conditional regularity]
\label{lem:head-mismatch-regular}
Under Assumption~\ref{ass:head-regularity-k1},
\begin{equation}
  M_{\Delta V}^\infty
  \le
  L_g \sqrt{B_N} + \mathcal R_N.
  \label{eq:head-mismatch-regular}
\end{equation}
\end{lemma}

\begin{proof}
For each coordinate $j\in\{1,2\}$,
\[
  |(v_i^A-v_i^B)_j|
  \le
  |(g(\xi_i^A)-g(\xi_i^B))_j|
  +
  |(r_i^A-r_i^B)_j|
  \le
  L_g\norm{\xi_i^A-\xi_i^B} + |(r_i^A-r_i^B)_j|.
\]
Take the empirical $\ell_2$ norm in $i$ and then the maximum over $j$.
\end{proof}

\begin{proposition}[Refined vector-head transfer]
\label{prop:refined-transfer}
Assume Assumption~\ref{ass:confinement-input}, and assume in addition that
$\nabla_\xi\varphi$ is Lipschitz in $\xi$ with constant $L_{\nabla}$ uniformly
in $x$. Then
\begin{equation}
  \sup_{t\in[0,1]}\ \sup_{x\in\mc{X}}
  \norm{z_t(x)-\hat z_t(x)}_\infty
  \le
  \frac{K_{\nabla}}{4} M_{\Delta V}^\infty \sqrt{B_N}
  +
  \frac{L_{\nabla}}{8} M_{V,T} B_N,
  \label{eq:refined-transfer}
\end{equation}
where $M_{V,T}:=\max\{A_T,C_T\}$.
\end{proposition}

\begin{proof}
Write
\[
  \phi_i^A(x):=\varphi(\xi_i^A,x),
  \qquad
  \phi_i^B(x):=\varphi(\xi_i^B,x),
  \qquad
  \phi_i^t(x):=\varphi((1-t)\xi_i^A+t\xi_i^B,x).
\]
Then
\begin{align*}
  z_t(x)-\hat z_t(x)
  &=
  \frac1N\sum_{i=1}^N
  \Bigl(
    v_i^t\phi_i^t-(1-t)v_i^A\phi_i^A-tv_i^B\phi_i^B
  \Bigr)
  \\
  &=
  \frac1N\sum_{i=1}^N
  v_i^t
  \Bigl(
    \phi_i^t-(1-t)\phi_i^A-t\phi_i^B
  \Bigr)
  \\
  &\qquad
  +
  t(1-t)\frac1N\sum_{i=1}^N
  (v_i^B-v_i^A)(\phi_i^A-\phi_i^B),
\end{align*}
where $v_i^t:=(1-t)v_i^A+tv_i^B$.

For the second sum, use
\[
  |\phi_i^A(x)-\phi_i^B(x)|
  \le
  K_{\nabla}\norm{\xi_i^A-\xi_i^B}
\]
to obtain, coordinate-wise,
\[
  t(1-t)\frac1N\sum_{i=1}^N
  |(v_i^B-v_i^A)_j|\,|\phi_i^A(x)-\phi_i^B(x)|
  \le
  \frac{K_{\nabla}}{4} M_{\Delta V}^\infty \sqrt{B_N}.
\]

For the first sum, the Lipschitz property of $\nabla_\xi\varphi$ gives the
second-order interpolation remainder
\[
  \bigl|
    \phi_i^t(x)-(1-t)\phi_i^A(x)-t\phi_i^B(x)
  \bigr|
  \le
  \frac{L_{\nabla}}{2}\,t(1-t)\norm{\xi_i^A-\xi_i^B}^2.
\]
Since $| (v_i^t)_j |\le M_{V,T}$ under
Assumption~\ref{ass:confinement-input},
\[
  \frac1N\sum_{i=1}^N
  |(v_i^t)_j|
  \bigl|
    \phi_i^t(x)-(1-t)\phi_i^A(x)-t\phi_i^B(x)
  \bigr|
  \le
  \frac{L_{\nabla}}{8} M_{V,T} B_N.
\]
Adding the two bounds and taking the maximum over $j\in\{1,2\}$ proves
\eqref{eq:refined-transfer}.
\end{proof}

\begin{proposition}[Endpoint raw-scale rate under head regularity]
\label{prop:scale-rate}
Assume Assumption~\ref{ass:confinement-input}. Then
\begin{equation}
  \Delta_{s,N}^{1/2}
  \le
  K_\varphi M_{\Delta V}^\infty + C_T K_{\nabla}\sqrt{B_N}.
  \label{eq:scale-rate}
\end{equation}
Consequently, if
\[
  M_{\Delta V}^\infty = O(\sqrt{B_N}),
\]
then
\[
  \Delta_{s,N} = O(B_N).
\]
\end{proposition}

\begin{proof}
Exactly as in the proof of Lemma~\ref{lem:w1-permutation},
\begin{align*}
  s_{\Theta_A}(x)-s_{\Theta_B}(x)
  &=
  \frac1N\sum_{i=1}^N
  \bigl(c_i^A-c_i^B\bigr)\varphi(\xi_i^A,x)
  \\
  &\qquad
  +
  \frac1N\sum_{i=1}^N
  c_i^B\bigl(\varphi(\xi_i^A,x)-\varphi(\xi_i^B,x)\bigr).
\end{align*}
The first term is bounded by
$K_\varphi M_{\Delta V}^\infty$ and the second by
$C_TK_{\nabla}\sqrt{B_N}$. Take the supremum in $x$ and then the
$L^2(P_X)$ norm.
\end{proof}

\begin{corollary}[Fast barrier under conditional head regularity]
\label{cor:fast-rate}
Assume Assumptions~\ref{ass:confinement-input}
and~\ref{ass:head-regularity-k1}. Assume also that
$\nabla_\xi\varphi$ is Lipschitz with constant $L_{\nabla}$ and that
\[
  \mathcal R_N = O(\sqrt{B_N}).
\]
Then
\[
  \cB(\Theta_A,\Theta_B;\pi)=O(B_N).
\]
In particular, if in addition
\[
  B_N = O(d/N),
\]
then
\[
  \cB(\Theta_A,\Theta_B;\pi)=O(d/N),
\]
which is $O(1/N)$ for fixed input dimension $d$.
\end{corollary}

\begin{proof}
Lemma~\ref{lem:head-mismatch-regular} implies
$M_{\Delta V}^\infty=O(\sqrt{B_N})$.
Proposition~\ref{prop:refined-transfer} then yields
\[
  \sup_t\sup_x \norm{z_t(x)-\hat z_t(x)}_\infty = O(B_N).
\]
Because $\Gamma$ is $1$-Lipschitz, the same $O(B_N)$ bound holds for
$\eta_{\Theta_t}-\tilde\eta_t$. Proposition~\ref{prop:scale-rate} gives
$\Delta_{s,N}=O(B_N)$, hence Proposition~\ref{prop:softplus-jensen} gives
\[
  \sup_t \E\big[\norm{\tilde\eta_t(X)-\hat\eta_t(X)}_\infty\big] = O(B_N).
\]
Repeating the proof of Theorem~\ref{thm:deterministic-lmc} with these refined
estimates gives $\cB=O(B_N)$. The final statement follows by substituting
$B_N=O(d/N)$.
\end{proof}

\section{From the Companion Mean-Field Theorem to LMC}

The companion note proves a training-side common-path theorem in the
continuous-time full-batch regime: if two width-$N$ runs are initialized i.i.d.\
from the same law and trained up to a fixed finite horizon $T$, then their
empirical parameter measures are both close to the same deterministic
mean-field law. The purpose of this section is to turn that statement into an
LMC corollary.

\begin{assumption}[Finite-horizon confinement input]
\label{ass:confinement-input}
Fix a training horizon $T<\infty$. Assume the trained endpoints satisfy
\[
  |a_i^A|,\ |a_i^B|\le A_T,
  \qquad
  |c_i^A|,\ |c_i^B|\le C_T,
  \qquad
  \norm{\xi_i^A},\ \norm{\xi_i^B}\le R_T
\]
for all $i=1,\dots,N$.
\end{assumption}

\begin{lemma}[Permutation selected by empirical Wasserstein distance]
\label{lem:w1-permutation}
Let
\[
  \rho_A^N:=\frac1N\sum_{i=1}^N \delta_{\theta_i^A},
  \qquad
  \rho_B^N:=\frac1N\sum_{i=1}^N \delta_{\theta_i^B},
  \qquad
  \theta_i^X:=(a_i^X,c_i^X,\xi_i^X).
\]
Let $\varepsilon_N:=W_1(\rho_A^N,\rho_B^N)$, where $W_1$ is computed with the
Euclidean cost on $\R^{2+p}$. Then there exists a permutation $\pi_N$ such that
\begin{equation}
  \frac1N\sum_{i=1}^N \norm{\theta_i^A-\theta_{\pi_N(i)}^B}
  =
  \varepsilon_N.
  \label{eq:w1-permutation}
\end{equation}
After relabeling network $B$ by this permutation, the quantities
$B_N$ and $\Delta_{s,N}$ satisfy
\begin{align}
  B_N
  &\le
  2R_T\,\varepsilon_N,
  \label{eq:w1-to-BN}
  \\
  \Delta_{s,N}
  &\le
  \bigl(K_\varphi + C_T K_{\nabla}\bigr)^2\,\varepsilon_N^2.
  \label{eq:w1-to-scale}
\end{align}
\end{lemma}

\begin{proof}
Because both empirical measures assign mass $1/N$ to each atom, the optimal
transport plan for $W_1$ is an optimal bistochastic matrix. By Birkhoff's
theorem, an extreme-point optimizer is a permutation, proving
\eqref{eq:w1-permutation}.

For \eqref{eq:w1-to-BN}, note that
\[
  \norm{\xi_i^A-\xi_i^B}\le 2R_T
  \qquad\text{and}\qquad
  \norm{\xi_i^A-\xi_i^B}\le \norm{\theta_i^A-\theta_i^B}.
\]
Hence
\[
  \norm{\xi_i^A-\xi_i^B}^2
  \le
  2R_T \norm{\xi_i^A-\xi_i^B}
  \le
  2R_T \norm{\theta_i^A-\theta_i^B}.
\]
Averaging gives \eqref{eq:w1-to-BN}.

For \eqref{eq:w1-to-scale}, write
\begin{align*}
  s_{\Theta_A}(x)-s_{\Theta_B}(x)
  &=
  \frac1N\sum_{i=1}^N
  \bigl(c_i^A-c_i^B\bigr)\varphi(\xi_i^A,x)
  \\
  &\qquad
  +
  \frac1N\sum_{i=1}^N
  c_i^B\bigl(\varphi(\xi_i^A,x)-\varphi(\xi_i^B,x)\bigr).
\end{align*}
Using $|\varphi|\le K_\varphi$, $|c_i^B|\le C_T$, and the gradient bound on
$\varphi$,
\[
  |s_{\Theta_A}(x)-s_{\Theta_B}(x)|
  \le
  K_\varphi \frac1N\sum_{i=1}^N |c_i^A-c_i^B|
  +
  C_T K_{\nabla}\frac1N\sum_{i=1}^N \norm{\xi_i^A-\xi_i^B}.
\]
Since both $|c_i^A-c_i^B|$ and $\norm{\xi_i^A-\xi_i^B}$ are bounded by
$\norm{\theta_i^A-\theta_i^B}$,
\[
  \sup_{x\in\mc{X}} |s_{\Theta_A}(x)-s_{\Theta_B}(x)|
  \le
  \bigl(K_\varphi + C_T K_{\nabla}\bigr)\varepsilon_N.
\]
Squaring and taking expectation in $X$ proves \eqref{eq:w1-to-scale}.
\end{proof}

\begin{corollary}[Finite-horizon LMC from a shared initialization law]
\label{cor:lmc-from-common-path}
Assume Assumption~\ref{ass:confinement-input}. Let $\rho_A^N$ and $\rho_B^N$
be the empirical measures of two independently trained width-$N$ runs of the
single-Gaussian natural-softplus network, both started from the same
initialization law and trained up to the same finite horizon $T$ in the
continuous-time full-batch regime of the companion note. Suppose that
\[
  \varepsilon_N:=W_1(\rho_A^N,\rho_B^N)\to 0
\]
almost surely, hence also in probability, as $N\to\infty$.
Let $\pi_N$ be the permutation from Lemma~\ref{lem:w1-permutation}. Then
\begin{equation}
  \cB(\Theta_A,\Theta_B;\pi_N)
  \le
  C_\eta(K_\varphi M_{V,T},\lambda_{\min},Y_\ast)
  \left[
    \frac12 K_{\nabla} M_{V,T}\sqrt{2R_T}\,\varepsilon_N^{1/2}
    +
    \frac{(K_\varphi + C_TK_{\nabla})^2}{32}\,\varepsilon_N^2
  \right],
  \label{eq:lmc-common-path-bound}
\end{equation}
where $M_{V,T}:=\max\{A_T,C_T\}$. In particular,
\[
  \cB(\Theta_A,\Theta_B;\pi_N)\to 0
\]
almost surely, hence also in probability, as $N\to\infty$.
\end{corollary}

\begin{proof}
Assumption~\ref{ass:confinement-input} implies $M_V^\infty\le M_{V,T}$.
Apply Theorem~\ref{thm:deterministic-lmc} after relabeling network $B$ by the
permutation $\pi_N$. Then use Lemma~\ref{lem:w1-permutation} to replace $B_N$
and $\Delta_{s,N}$ by the Wasserstein-based bounds \eqref{eq:w1-to-BN} and
\eqref{eq:w1-to-scale}. Since $\varepsilon_N\to 0$, the right-hand side of
\eqref{eq:lmc-common-path-bound} converges to zero.
\end{proof}

\begin{remark}[What the common-path corollary does and does not give]
Corollary~\ref{cor:lmc-from-common-path} is the clean Ferbach-style bridge from
training dynamics to connectivity, but it is not yet the sharpest rate
statement in this note. By itself it gives
\[
  \cB = O(\varepsilon_N^{1/2}+\varepsilon_N^2),
  \qquad
  \varepsilon_N=W_1(\rho_A^N,\rho_B^N).
\]
The faster $O(d/N)$ route of Corollary~\ref{cor:fast-rate} needs extra
structure: it uses hidden-weight alignment together with conditional head
regularity of the trained law.
\end{remark}

\begin{remark}[The rate-limiting term]
The bound \eqref{eq:lmc-common-path-bound} has two pieces:
\begin{enumerate}[label=(\alph*)]
  \item the transfer term, of order $\varepsilon_N^{1/2}$;
  \item the softplus Jensen term, of order $\varepsilon_N^2$.
\end{enumerate}
Thus, in this first $K=1$ theorem, the barrier is dominated by the
parameter-to-output transfer step. This cleanly separates the role of the
single-Gaussian convex output geometry from the role of neuron alignment.
\end{remark}

\section{What This Achieves}

This note deliberately stops before the MDN case. For $K=1$, the structure is:
\begin{enumerate}[label=(\roman*)]
  \item the training-side companion note supplies a finite-horizon common-path
  theorem;
  \item the present note turns that common path into an aligned barrier bound;
  \item output-space convexity is exact in the natural parameter $\eta$.
\end{enumerate}

This is the cleanest first LMC theorem available in the present project.
Compared with the older loss-regularity framework, the present note isolates the
precise reason the $K=1$ case is easier:
\begin{quote}
there is no component competition, so the loss interpolation is handled by
convexity in $\eta$, and the only remaining defects come from vector-head
transfer and softplus curvature.
\end{quote}

\end{document}
